\subsection{外代数的反交换性}

命题5. (i) 设 \((x_{i})_{1\leqslant i\leqslant n}\) 为模 \(\mathbf{M}\) 的元素组成的有限序列；对于每个置换 \(\sigma\) （对称群 \(\mathfrak{S}_n\) 中的），

\[
x _ {\sigma (1)} \wedge x _ {\sigma (2)} \wedge \dots \wedge x _ {\sigma (n)} = \varepsilon_ {\sigma}. x _ {1} \wedge x _ {2} \wedge \dots \wedge x _ {n}\tag{6}
\]

其中 \(\varepsilon_{\sigma}\) 表示置换 \(\sigma\) 的符号。

\begin{enumerate}
\def\labelenumi{(\roman{enumi})}
\setcounter{enumi}{1}
\tightlist
\item
  如果存在两个不同的指标 \(i, j\) 使得 \(x_{i} = x_{j}\) ，则乘积
\end{enumerate}

\[
x _ {1} \wedge x _ {2} \wedge \dots \wedge x _ {n}
\]

为零。

\begin{enumerate}
\def\labelenumi{(\roman{enumi})}
\tightlist
\item
  首先，由于 \(x \wedge x = 0\) 对所有 \(x \in \mathbf{M}\) 成立（由理想 \(\mathfrak{J}''\) 的定义可知），此外，对于 \(x, y\) 属于 \(\mathbf{M}\) 的情形，
\end{enumerate}

\[
x \wedge y + y \wedge x = (x + y) \wedge (x + y) - x \wedge x - y \wedge y = 0.
\]

这确立了（6）在 \(n = 2\) 的情形。一般情形随后由 §4 第 6 号的引理 3 得出。

\begin{enumerate}
\def\labelenumi{(\roman{enumi})}
\setcounter{enumi}{1}
\tightlist
\item
  在 (ii) 的假设下，存在一个置换 \(\sigma \in \mathfrak{S}_n\) ，使得 \(\sigma(1) = i\) 和 \(\sigma(2) = j\) ；那么对这个置换而言，(6) 的左边为零，因此右边也为零。
\end{enumerate}

推论 1. 设 H, K 为 N 的区间 \([1, n]\) 的两个互补子集，并设 \((i_{h})_{1 \leqslant h \leqslant p}, (j_{k})_{1 \leqslant k \leqslant n - p}\) 分别按递增顺序排列的 H 和 K 的元素序列；我们记

\[
x _ {H} = x _ {i _ {1}} \wedge x _ {i _ {2}} \wedge \dots \wedge x _ {i _ {p}}, \quad x _ {K} = x _ {j _ {1}} \wedge x _ {j _ {2}} \wedge \dots \wedge x _ {j _ {n - p}};
\]

那么

\[
x _ {\mathrm{H}} \wedge x _ {\mathrm{K}} = (- 1) ^ {\nu} x _ {1} \wedge x _ {2} \wedge \dots \wedge x _ {n}\tag{7}
\]

其中 \(\nu\) 是有序对 \((i,j)\in \mathbf{H}\times \mathbf{K}\) 中满足 \(i > j\) 者的个数

根据命题5，这归结为证明

引理 1. 若 \(\sigma \in \mathfrak{S}_n\) 是满足 \(\sigma(h) = i_h\) （对 \(1 \leqslant h \leqslant p\) ）与 \(\sigma(h) = j_{h-p}\) （对 \(p + 1 \leqslant h \leqslant n\) ）的置换，则 \(\varepsilon_\sigma = (-1)^{\nu}\) .

对于 \(1 \leqslant h < h' \leqslant p\) 或 \(p + 1 \leqslant h < h' \leqslant n\) ，有 \(\sigma(h') > \sigma(h)\) ，而有序对 \((h, h')\) ——即满足 \(1 \leqslant h \leqslant p < h' \leqslant n\) 与 \(\sigma(h) > \sigma(h')\) 者——的个数等于 \(\nu\) .

推论 2. 分次代数 \(\bigwedge (\mathbf{M})\) 是交错的（§ 4，第 9 号）。

只需将 §4 第 9 号的命题 13 应用于 \(\bigwedge (\mathbf{M})\) ，取生成系为集合 \(\mathbf{M}\) 并利用命题5。

命题6. 若 M 是有限生成的 A-模，则 \(\bigwedge (\mathbf{M})\) 是有限生成的 A-模；此外，若 M 具有含 \(n\) 个元素的生成系，则 \(\bigwedge^{p}(\mathbf{M}) = \{0\}\) 对于 \(p > n\) .

设 \((x_{i})_{1\leqslant i\leqslant n}\) 为 M 的一个生成系。\(\bigwedge^{p}(\mathbf{M})\) 中的每个元素都是形如

\[
x _ {i _ {1}} \wedge x _ {i _ {2}} \wedge \dots \wedge x _ {i _ {p}}
\]

其中指标 \(i_k\) 属于 \([1, n]\) ；根据命题 5，可设这些指标互不相同（否则相应的元素为零）。如果 \(p > n\) ，则不存在这样的指标序列，因此 \(\bigwedge^p(\mathbf{M}) = \{0\}\) 。若 \(p \leqslant n\) ，这样的序列只有有限多个，这就完成了证明。

